\begin{resumo}

%O resumo deve conter no máximo 500 palavras, devendo ser justificado na largura da página e escrito em um único parágrafo\footnote{E também não deve ter notas de rodapé; em outras palavras, não siga este exemplo... ;-)} com um afastamento de 1,27 cm na primeira linha. O espaçamento entre linhas deve ser de 1,5 linhas. O resumo deve ser informativo, ou seja, é a condensação do conteúdo e expõe finalidades, metodologia, resultados e conclusões.

%\lipsum[10-13]	% texto aleatório

A geração de conteúdo procedural é extremamente relevante e desafiadora na indústria dos videojogos, tópico que se tornou ainda mais atual com os modelos neurais profundos. Em especial, o aprendizado adversarial é muito utilizado para a geração de níveis de jogos 2D. Este trabalho utiliza um modelo neural adversarial generativo conhecido para gerar níveis do jogo \textit{Super Mario Bros}(1983). Para diversificar ainda mais os níveis gerados, foram utilizadas técnicas que impõem padrões para o modelo generativo seguir que tem como base algoritmos da computação evolutiva (tanto mono-objetivo quanto multi-objetivo), bem como conjuntos de treinamento diversos. Os algoritmos evolutivos foram comparados em desempenho, tempo de execução, adesão ao padrão do conjunto de treinamento e jogabilidade, avaliação esta última feita por um agente conceituado na literatura. Os resultados permitem gerar níveis diversos e jogáveis, além de manipular estes níveis para atenderem propriedades e dificuldades desejadas.

\end{resumo}
